\scanpage{22}
\textcircled{\mspage{15}}

On fixe entier $n$ premier aux car. résiduelles

\uline{Classe fondamentale locale}

Supposons $Y\longrightarrow X$ ``intersection complète de dim relative $1$'' et supposons d'abord $Y=V(\varphi)$, $\varphi$ section régulière de $\mathcal O_X$. Alors si $u:X-Y=U\to X$, donc \deleted{$\unsure{i|u}_*(\mathcal O_U)\sim\mathcal O_X[\frac1\varphi]$},
\deleted{$\unsure{i|u}_*(\mathcal O_U^*)$ \unsure{$\to$}}
\[
\varphi\in\Gamma(\unsure{i|u}_*(G_{mU}))
\]
\deleted{$\text{et }\unsure{n}\in\Gamma(X,\unsure{i|u}_*(G_{mU}^*))$}

\deleted{D'où} alors $\varphi$ est déterminé \inserted{par $Y$} mod multiplication par un $\in\Gamma(X,\unsure{i|u}_*(G_{mX}))$, donc \deleted{l'image} $\varphi$ définit une section \deleted{$\varphi\in$} \uline{canonique} de \deleted{$\unsure{i|u}_*(\mathcal O_U^*)/$} $\unsure{i|u}_*(G_{mU})/G_{mX}$. Or on a la suite exacte canonique
\diagram{\begin{tikzcd}[column sep=small]
0\arrow[r]&\underline H^0_Y(G_{mX})\arrow[r]\arrow[d,equal]&G_{mX}\arrow[r,no head]&i_*(G_{mU})\arrow[r,no head,"d"]&\underline H^1_Y(G_{mX})\arrow[d]\\
&0&&&0
\end{tikzcd}}
qui donne ici
\[
\deleted{\unsure{H^1}}\ \underline H^1_Y(G_{mX})\xrightarrow{\sim}i_*(G_{mU})/\operatorname{Im}G_{mX}\ ,
\]
et montre que $d\varphi$ est une section \uline{canonique} de $\underline H^1_Y(G_{mX})$, que nous allons noter \deleted{$\unsure{\varepsilon}_{Y/X}$} $\underline{\unsure{\beta}}_{Y/X}$. D'autre part, on a \uline{la} suite exacte de Kummer
\[
0\mathbin{-}\mu_{nX}\longrightarrow G_{mX}\mathrel{\overset{n}{-}}G_{mX}\longrightarrow0
\]
\mspage{21}
